% =====================================================================
% Draft article: Methods 1-2 only -- prompted and fine-tuned generative
% judges (direct and marker-conditioned) for RAG reliability.
% Build from this directory: pdflatex reliability_prompted_lora_judges_article.tex
% Run pdflatex twice for stable references.
% NOTE: zapiski.cls is not vendored in this repository; obtain it from
% github.com/madrugado/template-zapiski before building.
% =====================================================================
\documentclass{zapiski}
\originfo{529}{}{}{2026}
\setcounter{page}{1}
\date{}

\usepackage[cp1251]{inputenc}
\usepackage[T2A]{fontenc}
\usepackage[russian,english]{babel}
\usepackage{times}
\usepackage{url}
\usepackage{latexsym}
\usepackage{amssymb}
\usepackage{amsmath}
\usepackage{booktabs}
\usepackage{tabularx}
\usepackage{multirow}
\usepackage{caption}
\usepackage{cite}
\usepackage{float}

\newcommand{\fmacro}{\mathrm{F1}_{\mathrm{macro}}}
\renewcommand{\UrlFont}{\ttfamily\small}

\english

\begin{document}

\title[Prompted and Fine-Tuned RAG Judges]{Prompted and Fine-Tuned Generative
Judges for RAG Reliability: Axis Collapse, Error-Seeking Bias, and the Limits
of Marker Conditioning}

\author[Team~14\_1]{\text{Aleksandr Demchuk}}
\author[Team~14\_1]{\text{Murad Rezvan}}
\author[Team~14\_1]{\text{Sergey Natalenko}}
\author[Team~14\_1]{\text{Warda Tariq}}

\begin{abstract}
We study generative LLM-as-judge methods for automatic reliability assessment
of a Russian-language banking retrieval-augmented generation (RAG) system. A
response is reliable if it is both faithful to the retrieved chunks and
relevant to the client dialogue. Two prompt contracts are compared: a
\emph{direct} judge that emits the two binary labels as JSON, and a
\emph{marker-conditioned} judge that first names an error type from a
21-label taxonomy and then emits the same labels. Each contract is evaluated
zero-shot on Qwen2.5-1.5B-Instruct, after LoRA adaptation of the same model,
and after full supervised fine-tuning of Qwen2.5-7B-Instruct. On the canonical
1,796/224/225 split the always-reliable floor is 0.4186 reliability macro-F1.
The zero-shot direct judge reaches 0.4701, a difference of $+0.0515$ whose
paired bootstrap interval [$-0.0124$, $0.1140$] includes zero; the
marker-conditioned judge reaches 0.3513 and is the only variant tested here
that is \emph{significantly worse than the trivial baseline}
($\Delta=-0.0673$, [$-0.1190$, $-0.0182$], $p=0.004$), and the only method in
the project whose reported score falls below that baseline at all. We localize three
mechanisms. First, \emph{axis collapse}: the prompted judge emits identical
values on both axes in 99.87\% of rows while the gold labels disagree on
16.70\%, so the two-axis contract is not realized at all. Second, an
\emph{error-seeking bias} rather than the PASS bias reported for larger
judges: the direct judge flags 57.2\% of the corpus as unreliable against a
27.7\% base rate, and its PASS verdict leaves 23.3\% unreliable answers in the
approved set. Third, marker conditioning multiplies the error-seeking bias
(89.1\% flagged) while 27.8\% of its marker strings fall outside the closed
vocabulary that the prompt states, so the marker macro-F1 of 0.0123 confounds
taxonomy incompetence with format non-compliance. Balancing the supervised
training set is the only lever that helps: full fine-tuning of the 7B model
with a balanced train split reaches 0.586 reliability macro-F1, the best score
obtained on this split by any method in the project, while marker
fine-tuning stays at 0.485 with marker macro-F1 of 0.079.
\end{abstract}

\keywords{RAG evaluation \and hallucination detection \and LLM-as-judge
\and faithfulness \and answer relevance \and LoRA \and supervised
fine-tuning.}

\maketitle

% =====================================================================
\section{Introduction}
\label{sec:introduction}

A retrieval-augmented generation system conditions a language model on
documents fetched from a curated knowledge base. Retrieval reduces unsupported
generation but guarantees neither that every statement follows from the
evidence nor that the answer addresses what the client actually asked. In a
banking deployment both failure modes carry regulatory and reputational cost,
and the quality gate must run locally because the inputs contain client data.

We consider two binary axes. Faithfulness indicates that the answer is
supported by the retrieved chunks, without mixing facts across chunks and
without dropping essential conditions. Relevance indicates that the answer
addresses the latest client request in the dialogue. The final target is
\begin{equation}
 y_{\mathrm{reliable}} = y_{\mathrm{faithful}} \land y_{\mathrm{relevant}},
 \label{eq:reliability}
\end{equation}
and the primary metric is macro-averaged F1 on this derived label.

This article restricts the comparison to the \emph{generative judge} family:
methods in which a decoder-only instruction model reads
$(\mathrm{dialogue}, \mathrm{context}, \mathrm{answer})$ and writes the labels
as text. Encoder-based alternatives on the same data and the same split are
reported separately. Within the generative family we cross two independent
design choices:

\begin{itemize}
 \item \textbf{Prompt contract} --- \emph{direct} (Method 1), where the model
 returns only the two binary labels, versus \emph{marker} (Method 2), where
 the model first names an error type from the annotator taxonomy and then
 returns the labels, following the rationale-as-intermediate-target scheme
 of~\cite{hsieh2023distilling}. The hypothesis is that forcing an explicit
 error-type decision improves label quality and yields diagnosable failure
 categories for free.
 \item \textbf{Adaptation regime} --- zero-shot prompting, parameter-efficient
 LoRA adaptation~\cite{hu2021lora} of a 4-bit 1.5B model in the style
 of~\cite{dettmers2023qlora}, and full supervised fine-tuning of a 7B model.
\end{itemize}

Our contributions are:
\begin{enumerate}
 \item a held-out evaluation of both prompt contracts across all three
 adaptation regimes on the canonical organizer split, with paired stratified
 bootstrap intervals for every comparison that has per-row predictions;
 \item the identification of \emph{axis collapse} --- the prompted judge
 emits one verdict twice rather than two independent judgements --- and a
 quantification of how much of the two-axis contract it forfeits;
 \item a characterization of the prompted judge's failure as an
 \emph{error-seeking bias}, the opposite of the PASS bias usually reported for
 larger judges, together with the decision-theoretic consequence that its PASS
 verdict is close to uninformative;
 \item a decomposition of the marker-mode failure into a taxonomy component
 and a format-compliance component, with 27.8\% out-of-vocabulary marker
 strings, and the corresponding named fix;
 \item a documented negative result on LoRA at 1.5B, including a
 previously unreported train/inference distribution mismatch: 61.9\% of rows
 lose text to the SFT length caps while inference runs on untruncated text;
 \item confirmation that train-set balancing, not model scale or learning
 rate, is the lever that moves this task, with the 7B balanced direct judge
 reaching 0.586 reliability macro-F1.
\end{enumerate}

% =====================================================================
\section{Related Work}
\label{sec:related}

Industrial RAG-evaluation frameworks such as RAGAS~\cite{es2023ragas}
formalize faithfulness and answer relevance and delegate both to an external
LLM judge. The reliability of that delegation is itself contested:
\cite{zheng2023judging} document position, verbosity and self-enhancement
biases in LLM judges and show that judge agreement with human annotators
degrades sharply outside the domains the judge was aligned on.
Lynx~\cite{ravi2024lynx} demonstrates that a fine-tuned open judge can surpass
a frontier API judge on hallucination detection, which motivates our
supervised branches; its HaluBench evaluation also shows how strongly judge
quality depends on language and domain, which is the regime we are in.

Using an error type as an intermediate generation target follows distilling
step-by-step~\cite{hsieh2023distilling}, where rationales supervise a smaller
student in addition to the label. Our Method 2 is the cheapest possible
instance of that idea: the rationale is a single categorical token sequence
drawn from the annotator taxonomy. A recurring practical obstacle for such
schemes is that free-form decoding does not respect a closed label set;
schema-guided or grammar-constrained decoding~\cite{willard2023efficient}
removes exactly that failure, and our marker analysis quantifies how much of
the observed damage it would prevent.

On the supervised side, LoRA~\cite{hu2021lora} and its quantized
variant~\cite{dettmers2023qlora} make judge fine-tuning feasible on consumer
hardware, and we use the Qwen2.5 instruction models~\cite{qwen2024qwen25} as
the base for every generative variant. RAGTruth~\cite{niu2024ragtruth}
provides the reference point for what token-level hallucination supervision
looks like when it is available; in our corpus it is not, which is why the
marker channel is the only rationale signal we can train on and why its
density matters so much.

% =====================================================================
\section{Data and Evaluation Protocol}
\label{sec:data}

\subsection{Corpus and Split}

The organizer corpus contains 2,245 production RAG cases. Every case carries
the full client dialogue, up to eight retrieved knowledge chunks, the bot
answer, two binary labels, and an optional list of annotator error markers.
Reliability is not stored; it is derived by Eq.~\ref{eq:reliability}. We use
the repository's stratified seed-42 split
(\texttt{dataset.split\_samples}, 80/10/10) throughout, which yields
1,796 / 224 / 225 rows.

Table~\ref{tab:joint-labels} reports the \emph{joint} label distribution
rather than the two marginals, because the joint structure is what the
two-axis prediction contract has to reproduce and what the axis-collapse
analysis in Section~\ref{sec:results-collapse} depends on.

\begin{table}[H]
\centering
\caption{Joint distribution of the two gold axes on the canonical seed-42
split. $F$ denotes faithfulness and $R$ relevance; the reliable class is the
$(1,1)$ cell.}
\label{tab:joint-labels}
\begin{tabular}{lrrrrrr}
\toprule
Split & Rows & $(0,0)$ & $(0,1)$ & $(1,0)$ & $(1,1)$ & $F \neq R$ \\
\midrule
Train      & 1,796 & 194 & 279 & 25 & 1,298 & 304 \\
Validation &   224 &  27 &  34 &  1 &   162 &  35 \\
Test       &   225 &  27 &  35 &  1 &   162 &  36 \\
\midrule
Full       & 2,245 & 248 & 348 & 27 & 1,622 & 375 \\
\bottomrule
\end{tabular}
\end{table}

Two properties of Table~\ref{tab:joint-labels} govern everything that follows.

First, the axes genuinely disagree on 375 rows, 16.70\% of the corpus. A judge
that emits one verdict for both axes therefore forfeits, by construction, any
credit on one sixth of the data. This is the quantity that
Section~\ref{sec:results-collapse} measures against.

Second, the split is stratified on the \emph{reliable} label only, not on the
joint pair. The 27 rows in the rare $(F{=}1, R{=}0)$ cell consequently land
25 / 1 / 1 rather than the 21 / 3 / 3 that joint stratification would give.
Validation and test each contain exactly one faithful-but-irrelevant row, so
on those splits the reliable label differs from the faithfulness label in a
single row. Per-axis relevance metrics on validation and test are computed
against a nearly degenerate contingency and should be read as such; this is a
protocol limitation of the shared split, not of any method evaluated on it.

The always-reliable baseline is 0.4186 reliability macro-F1 on the 225-row
test split, because 162 of 225 rows are reliable. On the full 2,245-row corpus
the same trivial predictor scores 0.4194. The two floors differ because the
populations differ; comparisons that mix them are not meaningful, and
Section~\ref{sec:protocols} states which of our own numbers belong to which
population.

\subsection{Marker Annotations}
\label{sec:data-markers}

The taxonomy is a closed set of 13 organizer \texttt{reason\_*} labels; the
repository schema additionally admits six legacy compact error labels and the
two sentinels \texttt{none} and \texttt{unknown}, giving 21 admissible strings
in total. Marker annotation is sparse and its provenance needs care
(Table~\ref{tab:marker-provenance}), because two different counts of it
circulate in our own project documents.

\begin{table}[H]
\centering
\caption{Marker annotation provenance. Occurrences exceed annotated rows
because 32 rows carry more than one marker; the conversion pipeline keeps only
the first marker per row and overwrites the marker of any gold-reliable row
with \texttt{none}.}
\label{tab:marker-provenance}
\footnotesize
\begin{tabular}{lr}
\toprule
Quantity & Count \\
\midrule
Rows with at least one annotated marker (raw corpus)          & 158 \\
\quad of which gold-reliable, marker discarded by conversion  &  18 \\
\quad of which gold-unreliable, first marker retained         & 140 \\
Total marker occurrences across all positions                 & 200 \\
Rows carrying more than one marker                            &  32 \\
\midrule
Usable explicit-marker rows after conversion                  & 140 \\
\quad as a fraction of the 2,245-row corpus                   & 6.24\% \\
\quad as a fraction of the 623 unreliable rows                & 22.47\% \\
\midrule
Explicit-marker rows in train / validation / test             & 106 / 14 / 20 \\
Distinct marker classes in train / validation / test          & 12 / 6 / 9 \\
Taxonomy classes with at least one usable row                 & 12 of 13 \\
\bottomrule
\end{tabular}
\end{table}

Three consequences follow. The 158-row figure and the 22\% figure quoted in
our team's overview report describe different populations and are both
correct only under that reading: 158 rows are annotated, but only 140 survive
conversion, and it is the 140 that give 22.47\% of unreliable cases. The
18 discarded rows are an annotation inconsistency in the corpus --- an error
marker attached to a row that both binary labels call reliable --- silently
resolved in favour of the binary labels. Finally, the test split contains 20
explicit-marker rows spread over 9 classes, roughly two rows per class, so
marker macro-F1 on the test split is not estimable to any useful precision
and we report it only as a floor-level indicator.

\subsection{Protocols and Evidence Tiers}
\label{sec:protocols}

Two evaluation populations appear in this article and are never mixed within
a table.

\textbf{Canonical test split (225 rows).} The primary protocol. Every headline
number is on this population.

\textbf{Full corpus (2,245 rows).} Reported for zero-shot variants only, where
it is legitimate because no training occurs and there is no split to leak
across. It is the population our team's overview report quoted for Methods 1
and 2, and we reproduce it so those numbers remain traceable. It is
\emph{not} used for LoRA or fine-tuned variants.

Evidence quality is also not uniform across variants, and the tables label it
explicitly:

\begin{itemize}
 \item \textbf{measured} --- per-row predictions are on disk, so confusion
 matrices, per-class scores, per-marker breakdowns and paired bootstrap
 intervals are all computed directly. This covers both zero-shot variants.
 \item \textbf{derived} --- per-row predictions were not persisted, but the
 reported macro-F1 values together with the reported prediction counts
 determine the confusion matrix uniquely. This covers the LoRA variants; the
 reconstruction and its internal consistency check are in
 Appendix~\ref{sec:appendix-lora}. The LoRA adapter weights are no longer on
 disk, so these rows are permanently derived rather than merely
 not-yet-measured.
 \item \textbf{reported} --- only aggregate macro-F1 from the training run is
 available, rounded to three decimals. This covers the 7B full fine-tuning
 variants. These rows carry \emph{no} uncertainty estimate and no
 diagnostics, and no statistical claim in this article rests on them.
\end{itemize}

\subsection{Metrics and Uncertainty}

We report macro-F1 for faithfulness, relevance and the derived reliability
label. Generative judges produce hard verdicts, not probabilities, so no
threshold is fitted anywhere in this article --- an important asymmetry with
the encoder methods evaluated on the same split, and one we return to in
Section~\ref{sec:discussion}.

Sampling uncertainty is estimated with 10,000 paired stratified bootstrap
replicates (seed 42), resampling separately within the reliable and unreliable
strata so that every replicate preserves the original class counts. All
compared systems are evaluated on the same resampled indices. Because the
strata sizes are held fixed, the always-reliable baseline is constant at
0.4186 across replicates, which turns any comparison against it into a clean
one-sample test.

% =====================================================================
\section{Methods}
\label{sec:methods}

\subsection{Prediction Contract and Parsing}
\label{sec:contract}

Every variant implements one contract: given
$(\mathrm{dialogue}, \mathrm{context}, \mathrm{answer})$ it returns
$\hat{y}_{\mathrm{faithful}}, \hat{y}_{\mathrm{relevant}} \in \{0,1\}$, plus a
marker string in marker mode. Raw model text is converted by a three-stage
fallback: balanced-brace JSON extraction, then a field-level regular
expression, then a conservative default of $(0,0)$ with an
\texttt{invalid\_output} flag. The conservative default is deliberate --- a
judge that emits garbage must not be scored as if it had approved the answer.

The context is assembled by concatenating the non-empty retrieved chunks under
explicit \texttt{[CHUNK $i$]} headers, so chunk boundaries are visible to the
model. This matters for the \texttt{reason\_chunk\_fact\_mixup} error class,
which is only detectable if the boundaries survive into the prompt.

\subsection{Method 1: Direct Judge}
\label{sec:m1}

The direct prompt states the task, gives one-sentence definitions of both
axes, instructs the model to evaluate relevance against the latest client
request when the question field contains a full dialogue, and requires the
output
\begin{center}
\texttt{\{"faithfulness": 0 or 1, "relevance": 0 or 1\}}
\end{center}
and nothing else. Instructions are in English while the payload is Russian;
instruction-following on Qwen2.5 is more reliable in English, and the axis
definitions are the only part of the prompt that must be followed literally.

\subsection{Method 2: Marker-Conditioned Judge}
\label{sec:m2}

The marker prompt is identical except that it additionally lists the 21
admissible marker strings, instructs the model to choose exactly one,
reserves \texttt{none} for answers that are both faithful and relevant and
\texttt{unknown} for unreliable answers that fit no other class, and requires
\begin{center}
\texttt{\{"marker": "...", "faithfulness": 0 or 1, "relevance": 0 or 1\}}
\end{center}
with the marker field \emph{first}, so that the error-type decision is
generated before the binary labels and can condition them. Decoding is
unconstrained: the closed vocabulary is stated in the prompt but not enforced
at the sampler, which is the design decision that
Section~\ref{sec:results-markers} shows to be load-bearing.

\subsection{Zero-Shot Inference}
\label{sec:zeroshot}

Both prompts are run against
\texttt{mlx-community/Qwen2.5-1.5B-Instruct-4bit} through the MLX backend,
wrapped as a single user turn in the tokenizer chat template, decoded greedily
with a 64-token budget. No context truncation is applied at inference: on the
training split the median direct prompt is 8,356 characters and the longest is
19,112, and every prompt in the corpus fits the model context window. The same chat-template wrapper is used at
training time, so the train and inference prompt shapes agree.

\subsection{LoRA Adaptation of the 1.5B Model}
\label{sec:lora}

Supervised targets are the gold labels serialized in exactly the output format
the prompt requests, so the SFT objective and the inference contract coincide.
Loss is masked to completion tokens only (\texttt{-{}-mask-prompt}): the judge
prompt is roughly fifty times longer than the JSON completion, and without
masking the gradient optimizes prompt reconstruction instead of label
prediction.

Adapters are trained with \texttt{mlx\_lm.lora}: rank 8, scale 20, dropout 0,
last 16 layers, batch size 1 with 8 gradient-accumulation steps, learning rate
$1\cdot10^{-4}$, \texttt{max\_seq\_length} 4096, 3,592 iterations for two
epochs over the 1,796 training rows. Sequence length was reached by
elimination: 2,048 truncates the assistant target away on long rows and
training reports \texttt{nan} loss, while 8,192 exhausts Apple Metal memory
near iteration 40 at approximately 29\,GB peak.

To keep prompts inside 4,096 tokens the SFT record builder caps the dialogue
at 2,000 characters and the context at 5,000 characters. This cap is the
source of a distribution mismatch that we quantify in
Appendix~\ref{sec:appendix-truncation} and that has not been reported before:
the caps are applied when writing the training, validation and SFT test
records, but the held-out evaluation file is written from the
\emph{untruncated} samples, and the inference command consumes that
untruncated file. The adapter is therefore trained on a systematically shorter
input distribution than the one it is tested on.

A second configuration balances the SFT training set: the four
$(\mathrm{faithfulness}, \mathrm{relevance})$ label pairs are resampled
deterministically to 256 rows each, giving 1,024 training rows, with
validation and test untouched. This isolates the effect of the label prior
from every other factor.

\subsection{Full Fine-Tuning of the 7B Model}
\label{sec:fullft}

The same two prompt contracts are used to fully fine-tune
\texttt{Qwen2.5-7B-Instruct} with the TRL \texttt{SFTTrainer} on a single
A100-80GB: bf16, gradient checkpointing, 8-bit AdamW,
\texttt{assistant\_only\_loss} enabled (the transformers-side equivalent of
the MLX prompt masking), two epochs, maximum sequence length 2,048, learning
rate $1\cdot10^{-5}$. Four variants were trained: direct, direct with a
balanced train split, marker, and direct-balanced at learning rate
$3\cdot10^{-5}$. Balancing here oversamples the rare label pairs up to the
majority pair size rather than downsampling to a fixed 256, so it differs from
the 1.5B balanced run in construction while sharing its intent. Evaluation
uses greedy decoding and the same parser and metric code as every other
variant.

% =====================================================================
\section{Results}
\label{sec:results}

Table~\ref{tab:main-results} is the main result. All rows are on the canonical
225-row test split and are therefore directly comparable to each other and to
the encoder methods evaluated on the same split.

\begin{table}[H]
\centering
\caption{Generative judges on the canonical 225-row test split, reliability
macro-F1 as the primary metric. Evidence tier is defined in
Section~\ref{sec:protocols}: measured rows carry full diagnostics and
bootstrap intervals, derived rows carry reconstructed confusion matrices,
reported rows carry neither.}
\label{tab:main-results}
\footnotesize
\begin{tabularx}{\linewidth}{Xlrrr}
\toprule
Variant & Tier & Faith. & Rel. & Reliable \\
\midrule
Always-reliable floor                    & ---      & 0.4201 & 0.4668 & 0.4186 \\
\midrule
\multicolumn{5}{l}{\emph{Zero-shot, Qwen2.5-1.5B-Instruct-4bit}} \\
Method 1, direct                         & measured & 0.4739 & 0.4238 & 0.4701 \\
Method 2, marker                         & measured & 0.3477 & 0.2399 & 0.3513 \\
\midrule
\multicolumn{5}{l}{\emph{LoRA, Qwen2.5-1.5B-Instruct-4bit}} \\
Direct, unbalanced                       & derived  & 0.4201 & 0.4668 & 0.4186 \\
Direct, balanced $256{\times}4$          & derived  & 0.4669 & 0.4668 & 0.4636 \\
Marker                                   & derived  & 0.4201 & 0.4668 & 0.4186 \\
\midrule
\multicolumn{5}{l}{\emph{Full fine-tuning, Qwen2.5-7B-Instruct}} \\
Direct, unbalanced                       & reported & 0.544  & 0.467  & 0.540  \\
\textbf{Direct, balanced}                & reported & \textbf{0.580} & \textbf{0.555} & \textbf{0.586} \\
Direct, balanced, lr $3\cdot10^{-5}$     & reported & 0.533  & 0.457  & 0.532  \\
Marker                                   & reported & 0.488  & 0.500  & 0.485  \\
\bottomrule
\end{tabularx}
\end{table}

\subsection{Neither Prompted Judge Establishes a Gain over the Floor}
\label{sec:results-zeroshot}

The zero-shot direct judge scores 0.4701 against a floor of 0.4186. The
paired stratified bootstrap puts that difference at $+0.0515$ with a 95\%
interval of [$-0.0124$, $0.1140$] and $P(\Delta>0)=0.945$. On the canonical
split the improvement is therefore \emph{not} established at the 95\% level.
This corrects the reading invited by our team's overview report, which places
the full-corpus figures 0.4946 and 0.4194 side by side without an interval;
the point estimate is positive and reasonably stable across populations
(0.4701 test, 0.4882 validation, 0.4946 full corpus), but 225 rows do not
resolve a gap of this size.

The marker judge scores 0.3513, and here the bootstrap is decisive in the
other direction: $\Delta=-0.0673$ against the floor with interval
[$-0.1190$, $-0.0182$] and $P(\Delta>0)=0.004$. Method 2 is the only variant
in this article --- and, on point estimates, the only method in the project
with reported numbers --- that is significantly \emph{worse than predicting a
constant}. We bootstrapped only our own variants against the floor, so the
comparison outside this article rests on point estimates alone.
The direct-versus-marker comparison is also decisive,
$P(\Delta>0)=0.9995$. Full intervals are in
Appendix~\ref{sec:appendix-bootstrap}.

Point estimates with intervals are in Table~\ref{tab:bootstrap-point} and the
paired comparisons in Table~\ref{tab:bootstrap-delta}; the underlying
confusion matrices for every split are in Tables~\ref{tab:direct-confusion}
and~\ref{tab:marker-confusion}.

Output-format compliance is not the explanation for either result. The direct
judge produces zero unparseable outputs on all 2,245 rows; the marker judge
produces three, an invalid rate of 0.13\%.

\subsection{Axis Collapse}
\label{sec:results-collapse}

The prompted judge does not perform two judgements. On the full corpus the
direct judge emits different values for the two axes on 3 rows out of 2,245
(0.13\%) and the marker judge on 2 rows (0.09\%), while the gold labels differ
on 375 rows (16.70\%); see Table~\ref{tab:axis-coupling}. The joint prediction
distribution (Table~\ref{tab:joint-predictions}) is essentially two-valued:
for the direct judge, 1,281 rows receive $(0,0)$, 961 receive $(1,1)$, and 3
receive $(0,1)$; the $(1,0)$ cell is never produced.

The rare gold cell makes this concrete. Of the 27 corpus rows that are
faithful but irrelevant, the direct judge labels 18 as $(0,0)$ and 9 as
$(1,1)$ --- never $(1,0)$, the correct answer. The marker judge labels all 27
as $(0,0)$.

The prompt states the axes as independent criteria and the model is not
capacity-bound on the format, since it emits well-formed JSON on every row.
What collapses is the decision, not the syntax: the judge forms a single
overall impression of the answer and writes it into both fields. Because
reliability is the conjunction, this collapse is partly masked in the headline
metric --- if both axes carry the same value, the conjunction equals that
value --- which is precisely why per-axis reporting is necessary to see it.
The practical consequence is that the whole two-axis design, including the
per-axis thresholds and the axis-independence instruction, buys nothing at
this model scale.

\subsection{Error-Seeking Bias, Not PASS Bias}
\label{sec:results-bias}

Large prompted judges in this project fail by approving too much. The small
prompted judge fails the opposite way. The direct judge flags 1,284 of 2,245
rows as unreliable, 57.2\%, against a gold base rate of 27.7\%. The marker
judge flags 2,000 rows, 89.1\%. Recall of the unreliable class is
correspondingly high --- 0.640 for direct and 0.905 for marker on the full
corpus --- and precision correspondingly poor, 0.311 and 0.282.

The decision-theoretic reading (Table~\ref{tab:operating-point}) is more
informative than the macro-F1. On the
full corpus the direct judge's PASS set contains 961 answers of which 23.3\%
are actually unreliable, against a 27.7\% base rate: accepting its approval
reduces the error density in the approved pile by four percentage points. On
the test split the same quantity is 27.0\% against a 28.0\% base rate, i.e.
no reduction at all. The PASS verdict of a 1.5B prompted judge is close to
uninformative, and a reviewer who trusts it is doing approximately as well as
a reviewer who skips reading.

Discrimination, as opposed to operating point, is uniformly weak. Youden's
$J = \mathrm{TPR}+\mathrm{TNR}-1$ on the reliability axis is $0.095$ (direct)
and $0.020$ (marker) on the full corpus, $0.022$ and $0.094$ on test, and
$0.080$ and $0.058$ on validation. The ordering between the two modes flips
between populations, so the large macro-F1 gap between direct and marker is
driven almost entirely by \emph{where each mode sits on the FAIL--PASS axis},
not by how well either separates the classes. Since a generative judge emits
a hard verdict, that operating point cannot be moved after the fact.

Per-marker recall (Appendix~\ref{sec:appendix-permarker}) is consistent with
this: the direct judge's recall on unreliable rows varies from 0.333
(\texttt{reason\_incomplete\_answer},
\texttt{reason\_missed\_chunk\_conditions}) to 0.783
(\texttt{reason\_answer\_for\_operator}), but the spread is dominated by class
sizes of 3--32 rows and is not separable from the 57\% blanket flag rate.

Prompt format matters more than any of this: the two contracts agree on only
60.6\% of reliability verdicts, Cohen's $\kappa = 0.113$
(Appendix~\ref{sec:appendix-agreement}).

Answer length is not the hidden driver. The AUC of answer length as a
predictor of the direct judge's unreliable verdict is 0.568, against 0.497 for
the gold unreliable label: the judge has a mild length bias that the labels
do not share, but it is far too weak to explain the flag rate.

\subsection{Marker Conditioning Fails in Two Separable Ways}
\label{sec:results-markers}

Marker macro-F1 is 0.0123 on the full corpus, which invites the summary
``noise''. That summary hides a decomposable structure.

\textbf{Format non-compliance.} 624 of 2,245 marker predictions, 27.8\%, are
strings outside the 21-label vocabulary the prompt states
(Table~\ref{tab:marker-vocab}). They are not
arbitrary: 566 are \texttt{reason\_hallucination} and 58 are
\texttt{reason\_hallucinated}, both one-token drifts from the admissible
\texttt{reason\_hallucinated\_fact}. The model has understood the naming
convention and failed to reproduce a specific string. Because the repository
metric averages over the union of gold-present and prediction-present labels,
each invented string enters as a guaranteed-zero class and drags the macro
average down. Restricting the average to gold-present classes raises the score
from 0.0123 to 0.0175 on the full corpus and from 0.0136 to 0.0222 on test
(Table~\ref{tab:marker-decomposition}).

\textbf{Taxonomy incompetence.} The remaining score is still floor-level, and
this component is real. Restricted to the explicit marker classes --- dropping
the \texttt{none} and \texttt{unknown} sentinels, which the metric otherwise
rewards for free --- macro-F1 is 0.0052 on the full corpus and 0.0000 on test.
Exact marker match on the 140 explicit-marker rows is 6, and on the 20 such
test rows it is 0. Of the 32 rows whose gold marker is
\texttt{reason\_hallucinated\_fact}, 13 receive a prediction somewhere in the
hallucination family, most of them the out-of-vocabulary spelling.

The two components have different fixes and different prognoses. Schema-guided
decoding~\cite{willard2023efficient} eliminates the first component outright
and would recover roughly a quarter of the predictions into the admissible
set. The second component is a data problem: 140 usable rows spread over 12 of the 13 taxonomy
classes --- \texttt{reason\_other} never survives conversion at all --- with
20 rows in the test split, cannot train or evaluate a 13-way classifier at any
decoding discipline.

Marker conditioning also \emph{costs} on the binary task it was supposed to
help. Requiring an error name first acts as an error-seeking prior: the model
that has just written \texttt{reason\_hallucination} is not going to write
\texttt{"faithfulness": 1}. This is the mechanism behind the 89.1\% flag rate
and, through it, behind the below-floor macro-F1. The hypothesis that markers
act as a free rationale is therefore rejected in its cheap form; whether it
survives with constrained decoding and a denser marker channel is untested.

\subsection{LoRA at 1.5B Does Not Beat Zero-Shot}
\label{sec:results-lora}

Unbalanced LoRA collapses completely: the adapter emits $(1,1)$ for all 225
held-out rows, reproducing the always-reliable floor exactly on every axis
(0.4201 / 0.4668 / 0.4186). Marker LoRA collapses identically on the binary
axes and additionally predicts \texttt{none} for every row, giving marker
macro-F1 of 0.0761. Two epochs of SFT on a corpus that is 72.3\% reliable
teaches the majority label and nothing else.

Balancing the SFT train set to 256 rows per label pair removes the
faithfulness collapse. The balanced adapter predicts $(1,1)$ on 133 rows and
$(0,1)$ on 92, lifting reliability macro-F1 from 0.4186 to 0.4636 and
faithfulness from 0.4201 to 0.4669; the reconstructed confusion matrices are
in Table~\ref{tab:lora-reconstruction}. It still never predicts
$\mathrm{relevance}=0$, so relevance macro-F1 stays pinned at the
all-positive value, and it still does not reach the zero-shot direct judge's
0.4701. Balancing recovers most of the damage that fine-tuning caused; it does
not produce a net gain over not fine-tuning at all.

Three factors plausibly contribute and we can only bound the first two. The
adapter is rank 8 over the last 16 layers of a 4-bit model, which is a small
capacity budget. The effective optimizer step count is 449, since the 3,592
iterations run at batch size 1 with 8 gradient-accumulation steps. And the
training distribution does not match the evaluation distribution: under the
2,000 / 5,000-character SFT caps, 61.9\% of rows lose text and 36.1\% of all
context characters are discarded, while the held-out evaluation file is
written untruncated (Appendix~\ref{sec:appendix-truncation}). The adapter
learns to judge from a truncated context and is then asked to judge from a
full one. This is a defect of the harness, not a property of LoRA, and it must
be fixed before the LoRA branch can be called properly tested.

\subsection{Full Fine-Tuning: Balancing Is the Lever}
\label{sec:results-fullft}

Scaling the same two contracts to a fully fine-tuned 7B model reproduces every
qualitative finding and shifts the level. Unbalanced direct fine-tuning
reaches 0.540 reliability macro-F1 but retains the relevance collapse ---
its relevance macro-F1 of 0.467 is the all-positive value of 0.4668 to
rounding, and the reported relevance confusion has no true negatives.
Balancing the train split lifts every metric at once: reliability
$0.540 \to 0.586$, faithfulness $0.544 \to 0.580$, relevance
$0.467 \to 0.555$. Raising the learning rate to $3\cdot10^{-5}$ undoes the
gain and lands the model \emph{below} the all-positive relevance value of
0.4668, at 0.457: with no true negatives reported, the handful of
$\mathrm{relevance}=0$ predictions it does emit are essentially all wrong, so
the high-LR model is strictly worse than collapsing. This confirms that the
lever was the label prior and not optimization strength. Marker fine-tuning again
underperforms direct on both binary axes (0.485 reliability) and leaves marker
macro-F1 at 0.079, the same floor as every other marker attempt.

At 0.586, the balanced 7B direct judge is the best score obtained on the
canonical 225-row test split by any method in this project, generative or
encoder-based. The encoder-based methods scored on this same split by our team
reach 0.5203 (validation-calibrated LettuceDetect features) and 0.5635
(pseudo-label-tuned claim-level NLI combined with a separate relevance head);
their split reproduces ours exactly on every count this comparison uses.
This claim needs one qualification that our own project documents get wrong. The RuModernBERT cross-encoder baseline frequently quoted
at 0.5879 is evaluated on a 1,571 / 225 / 449 split, not on the canonical
1,796 / 224 / 225 split, and its 449-row test population is twice the size and
differently constituted. The two numbers are not comparable, and the sentence
``full fine-tuning beats the encoder baseline'' is not supported by evidence
of the kind this article requires. What is supported is the within-split
statement above.

The 7B rows carry no uncertainty estimate. Given that the difference between
the balanced 7B judge and the encoder-family results on the same split is of
the order of 0.02--0.05 macro-F1, and that comparably sized differences in
Section~\ref{sec:results-zeroshot} did not survive a bootstrap on 225 rows, no
ranking among the leading methods should be treated as established.

% =====================================================================
\section{Discussion}
\label{sec:discussion}

\textbf{One verdict, written twice.} The single most consequential finding is
that the prompted 1.5B judge does not implement the two-axis contract at all.
It reads the answer, forms one impression, and writes it into both JSON
fields. Everything downstream inherits this: the axis-independence instruction
is inert, per-axis analysis is degenerate, and the 16.70\% of rows where the
axes genuinely disagree are unreachable by construction. Full fine-tuning at
7B partially breaks the collapse --- the balanced model's relevance macro-F1
of 0.555 is well above the all-positive value, so it is producing genuine
relevance decisions --- which suggests the collapse is a capacity and
supervision effect rather than a prompt-design effect. The design implication
is that if two axes are wanted from a small judge, they should be two calls or
two heads, not two fields of one JSON object.

\textbf{Operating point versus discrimination.} Both prompted variants have
near-zero discrimination ($J$ between 0.02 and 0.10, unstable across splits)
but very different operating points, and the entire macro-F1 gap between them
comes from the operating point. A generative judge that emits a hard verdict
gives no way to move that point after the fact. This is the structural
disadvantage of Methods 1 and 2 relative to methods that expose a score, and
it is worth stating precisely: the encoder methods on this same split gain
0.10 macro-F1 purely from validation-fitted thresholds, a lever that is simply
unavailable here. Reading verdict probabilities from token log-probabilities
--- the approach taken by the prompt-optimized judge branch of this project
--- is the natural way to recover it.

\textbf{Marker conditioning is not a free rationale.} The hypothesis was that
naming the error type first would sharpen the binary decision. It does the
opposite, and the mechanism is legible: committing to an error name commits
the model to finding an error. The cost is large enough to push the method
below the trivial baseline. The two failure components separate cleanly ---
one quarter of predictions violate a closed vocabulary that the prompt
states, and the rest are wrong about the taxonomy --- and only the first has a
cheap fix. At 140 usable annotated rows over 12 populated classes, the second
is not addressable at all on the current release.

\textbf{Prompting and fine-tuning fail from opposite ends of the same axis.}
The two failure modes in this article are one mechanism seen twice. Prompting
pushes the judge toward FAIL --- 57.2\% of the corpus flagged in direct mode
and 89.1\% in marker mode, against a 27.7\% base rate --- because the
instruction frames the task as fault-finding and the marker field turns that
framing into a commitment. Naive supervised fine-tuning pushes it the other
way, toward PASS, because two epochs on a 72.3\%-reliable corpus absorb the
label prior: the unbalanced adapter at 1.5B and the unbalanced full
fine-tune at 7B both report zero true negatives on relevance. Neither regime
is discriminating better than the other; they sit at opposite ends of the
FAIL--PASS axis. Balancing is the intervention that moves the supervised end
of that axis back toward the data, which is why it is the only lever that
worked at either scale.

\textbf{Balancing dominates scale and learning rate.} Across both model sizes,
the same single intervention moves the metric: making the four label pairs
equiprobable in the SFT set. It converts total collapse into a working
classifier at 1.5B and adds 0.046 macro-F1 at 7B, while a threefold learning
rate increase at 7B loses 0.054. The uniform failure mode of naive SFT here is
prior absorption, and the corpus prior is strong enough (72.3\% reliable,
87.7\% relevant) that two epochs of unbalanced training reach it and stop.

\textbf{The ceiling looks like a data property.} Every method in this
project, generative or encoder-based, lands between roughly 0.52 and 0.59
reliability macro-F1 on the canonical split, and the differences among the
leaders are of the same magnitude as the bootstrap intervals on 225 rows.
Going from a rank-8 adapter over 16 layers of a 1.5B model --- under a million
trainable parameters --- to all 7.6 billion parameters of the 7B model buys
about 0.12 macro-F1, while balancing a training set buys 0.05 at the same
scale. That pattern points at label quality, annotation density and test set
size rather than model capacity.

\textbf{Threats to validity.} The 225-row test split is small enough that no
comparison in the leading group is resolvable; the split is stratified on the
reliable label only, so per-axis relevance metrics on validation and test rest
on a single faithful-but-irrelevant row each; the 7B numbers are aggregates
without diagnostics or intervals; the LoRA numbers are reconstructions from
reported aggregates and cannot be re-measured because the adapter weights are
gone; and the LoRA branch is confounded by the SFT truncation mismatch, so its
negative result bounds the harness as configured rather than the method.

% =====================================================================
\section{Conclusion}
\label{sec:conclusion}

Within the generative-judge family, the prompt contract matters more than the
model, and the training-set prior matters more than both. The direct contract
beats the marker contract decisively at every scale tested. The
marker-conditioned judge is the only variant we tested against the trivial
baseline that performs significantly worse than a constant prediction, and it
is the only method in the project whose reported score falls below that
baseline at all; its failure decomposes
into a format-compliance component with a known fix and a data-density
component without one. Zero-shot prompting of a 1.5B judge does not establish
a gain over the trivial baseline on the canonical split, and its verdicts are
better described as error-seeking than as PASS-biased: it flags twice as many
answers as are actually unreliable, and its approvals carry almost no
information. LoRA at 1.5B does not beat its own zero-shot baseline, though
that comparison is confounded by a train/inference truncation mismatch that
should be repaired before the branch is written off. Full fine-tuning at 7B
with a balanced train split reaches 0.586 reliability macro-F1, the best score
on the canonical split in this project --- but without an uncertainty
estimate, and with the frequently quoted encoder comparison invalidated
because that baseline was scored on a different split.

The immediate work items follow directly: persist per-row predictions from the
7B evaluation so the leading result acquires the diagnostics and intervals
every other measured row in this article carries; apply schema-guided decoding
to the marker field; fix the SFT truncation mismatch and re-run the LoRA
branch; and, for any generative judge intended to be operated rather than
benchmarked, expose verdict probabilities so that the operating point becomes
a tunable parameter instead of a property of the prompt.

\medskip
\noindent\textbf{Code and artifacts.}
\url{https://github.com/aldem2k00/rag-reliability}

\clearpage
\appendix
\section{Detailed Results}
\label{sec:appendix-results}

\subsection{Zero-Shot Confusion Matrices and Per-Class Scores}
\label{sec:appendix-confusion}

\begin{table}[H]
\centering
\caption{Zero-shot direct judge (Method 1), Qwen2.5-1.5B-Instruct-4bit. TN,
FP, FN and TP are stated with respect to the positive class named in the axis
column: faithful, relevant, or reliable.}
\label{tab:direct-confusion}
\footnotesize
\begin{tabular}{llrrrrrrr}
\toprule
Split & Axis & TN & FP & FN & TP & F1$_0$ & F1$_1$ & $\fmacro$ \\
\midrule
\multirow{3}{*}{Full (2,245)}
 & Faithfulness & 381 & 215 & 903 & 746 & 0.4053 & 0.5716 & 0.4885 \\
 & Relevance    & 188 &  87 & 1{,}093 & 877 & 0.2416 & 0.5978 & 0.4197 \\
 & Reliability  & 399 & 224 & 885 & 737 & 0.4185 & 0.5707 & 0.4946 \\
\midrule
\multirow{3}{*}{Validation (224)}
 & Faithfulness &  38 &  23 &  90 &  73 & 0.4021 & 0.5637 & 0.4829 \\
 & Relevance    &  18 &  10 & 109 &  87 & 0.2323 & 0.5939 & 0.4131 \\
 & Reliability  &  39 &  23 &  89 &  73 & 0.4105 & 0.5659 & 0.4882 \\
\midrule
\multirow{3}{*}{Test (225)}
 & Faithfulness &  36 &  26 &  89 &  74 & 0.3850 & 0.5627 & 0.4739 \\
 & Relevance    &  18 &  10 & 106 &  91 & 0.2368 & 0.6107 & 0.4238 \\
 & Reliability  &  36 &  27 &  89 &  73 & 0.3830 & 0.5573 & 0.4701 \\
\bottomrule
\end{tabular}
\end{table}

\begin{table}[H]
\centering
\caption{Zero-shot marker judge (Method 2), same model and decoding. The
consistently large FN column is the error-seeking bias: the judge assigns the
negative class to almost everything.}
\label{tab:marker-confusion}
\footnotesize
\begin{tabular}{llrrrrrrr}
\toprule
Split & Axis & TN & FP & FN & TP & F1$_0$ & F1$_1$ & $\fmacro$ \\
\midrule
\multirow{3}{*}{Full (2,245)}
 & Faithfulness & 537 & 59 & 1{,}463 & 186 & 0.4137 & 0.1964 & 0.3051 \\
 & Relevance    & 253 & 22 & 1{,}745 & 225 & 0.2226 & 0.2030 & 0.2128 \\
 & Reliability  & 564 & 59 & 1{,}436 & 186 & 0.4300 & 0.1993 & 0.3146 \\
\midrule
\multirow{3}{*}{Validation (224)}
 & Faithfulness &  55 &  6 & 138 &  25 & 0.4331 & 0.2577 & 0.3454 \\
 & Relevance    &  25 &  3 & 168 &  28 & 0.2262 & 0.2467 & 0.2365 \\
 & Reliability  &  56 &  6 & 137 &  25 & 0.4392 & 0.2591 & 0.3491 \\
\midrule
\multirow{3}{*}{Test (225)}
 & Faithfulness &  59 &  3 & 140 &  23 & 0.4521 & 0.2434 & 0.3477 \\
 & Relevance    &  28 &  0 & 171 &  26 & 0.2467 & 0.2332 & 0.2399 \\
 & Reliability  &  60 &  3 & 139 &  23 & 0.4580 & 0.2447 & 0.3513 \\
\bottomrule
\end{tabular}
\end{table}

\begin{table}[H]
\centering
\caption{Operating point of the reliability verdict. The PASS set is the set
of answers the judge approves; leakage is the fraction of it that is actually
unreliable, to be compared with the base rate in the same population.
Youden's $J$ summarizes discrimination independently of operating point.}
\label{tab:operating-point}
\footnotesize
\begin{tabular}{llrrrrrr}
\toprule
Mode & Split & Base rate & PASS set & Leakage & FAIL set & FAIL prec. & $J$ \\
\midrule
\multirow{3}{*}{Direct}
 & Validation & 0.2768 &  96 & 0.2396 &  128 & 0.3047 & $+0.0796$ \\
 & Test       & 0.2800 & 100 & 0.2700 &  125 & 0.2880 & $+0.0220$ \\
 & Full       & 0.2775 & 961 & 0.2331 & 1{,}284 & 0.3107 & $+0.0948$ \\
\midrule
\multirow{3}{*}{Marker}
 & Validation & 0.2768 &  31 & 0.1935 &  193 & 0.2902 & $+0.0575$ \\
 & Test       & 0.2800 &  26 & 0.1154 &  199 & 0.3015 & $+0.0944$ \\
 & Full       & 0.2775 & 245 & 0.2408 & 2{,}000 & 0.2820 & $+0.0200$ \\
\bottomrule
\end{tabular}
\end{table}

The $J$ column is the reason Section~\ref{sec:results-bias} refuses to claim
that either mode discriminates better than the other: the ordering reverses
between test and full corpus, and every value is close to the zero of an
uninformative predictor. The leakage column is more stable and more
interpretable --- the direct judge's approvals are barely cleaner than the
population it drew them from.

\subsection{Axis Coupling}
\label{sec:appendix-coupling}

\begin{table}[H]
\centering
\caption{Rows on which the two predicted axes differ, against rows on which
the two gold axes differ. A judge implementing the two-axis contract should
approximately match the gold row.}
\label{tab:axis-coupling}
\begin{tabular}{lrrr}
\toprule
Source & Full (2,245) & Test (225) & Fraction (full) \\
\midrule
Gold labels                 & 375 & 36 & 0.1670 \\
Zero-shot direct prediction &   3 &  1 & 0.0013 \\
Zero-shot marker prediction &   2 &  0 & 0.0009 \\
\bottomrule
\end{tabular}
\end{table}

\begin{table}[H]
\centering
\caption{Joint prediction distribution against the joint gold distribution,
full corpus. The $(1,0)$ cell --- faithful but irrelevant --- is never
predicted by either mode.}
\label{tab:joint-predictions}
\footnotesize
\begin{tabular}{lrrrr}
\toprule
Source & $(0,0)$ & $(0,1)$ & $(1,0)$ & $(1,1)$ \\
\midrule
Gold                        &   248 & 348 & 27 & 1{,}622 \\
Zero-shot direct prediction & 1{,}281 &   3 &  0 &   961 \\
Zero-shot marker prediction & 1{,}998 &   2 &  0 &   245 \\
\bottomrule
\end{tabular}
\end{table}

On the 27 gold $(1,0)$ rows specifically, the direct judge predicts $(0,0)$
18 times and $(1,1)$ 9 times; the marker judge predicts $(0,0)$ 27 times.

\subsection{Prompt-Format Sensitivity}
\label{sec:appendix-agreement}

The two prompt contracts differ only in whether a marker field precedes the
binary fields. Their reliability verdicts agree on 60.6\% of the full corpus,
with Cohen's $\kappa = 0.113$ (test: 60.0\%, $\kappa = 0.125$). A minor
formatting change to the output contract therefore flips two rows in five.
This is a larger effect than the difference between any two adaptation
regimes we tested, and it bounds how much confidence any single-prompt
zero-shot result deserves.

\subsection{Bootstrap Uncertainty}
\label{sec:appendix-bootstrap}

Intervals are percentile 95\% intervals from 10,000 paired stratified
bootstrap replicates on the 225-row test split, seed 42, resampling within the
reliable and unreliable strata so class counts are preserved. All systems are
scored on identical resampled indices.

\begin{table}[H]
\centering
\caption{Test macro-F1 with bootstrap intervals, zero-shot variants. The
always-reliable baseline is constant across replicates because the strata
sizes are fixed.}
\label{tab:bootstrap-point}
\begin{tabular}{llr}
\toprule
Variant & Axis & $\fmacro$ [95\% CI] \\
\midrule
Always reliable  & Reliability  & 0.4186 [0.4186, 0.4186] \\
\midrule
\multirow{3}{*}{Direct}
 & Faithfulness & 0.4739 [0.4095, 0.5365] \\
 & Relevance    & 0.4238 [0.3642, 0.4834] \\
 & Reliability  & 0.4701 [0.4062, 0.5326] \\
\midrule
\multirow{3}{*}{Marker}
 & Faithfulness & 0.3477 [0.2963, 0.3982] \\
 & Relevance    & 0.2399 [0.1864, 0.2931] \\
 & Reliability  & 0.3513 [0.2996, 0.4004] \\
\bottomrule
\end{tabular}
\end{table}

\begin{table}[H]
\centering
\caption{Paired bootstrap comparisons on the reliability axis, test split.
$\Delta$ is the difference of the two point estimates; the interval and
$P(\Delta>0)$ come from the paired replicate distribution, in which the first
system scores above the second in that fraction of replicates.}
\label{tab:bootstrap-delta}
\footnotesize
\begin{tabular}{lrrr}
\toprule
Comparison & $\Delta$ & 95\% CI & $P(\Delta>0)$ \\
\midrule
Direct $-$ always reliable & $+0.0515$ & [$-0.0124$, $+0.1140$] & 0.9451 \\
Marker $-$ always reliable & $-0.0673$ & [$-0.1190$, $-0.0182$] & 0.0042 \\
Direct $-$ marker          & $+0.1188$ & [$+0.0451$, $+0.1929$] & 0.9995 \\
\bottomrule
\end{tabular}
\end{table}

The replicate means differ from the point differences by at most 0.0008
($-0.0681$ and $+0.1196$ for the second and third rows), the usual small
resampling bias; the point differences are reported because they are the
estimand.

Only two of the three comparisons are established at the 95\% level, and one
of them is a negative result: the marker judge is worse than a constant
prediction. The direct judge's advantage over the floor is not established on
225 rows.

\subsection{Per-Marker Recall of the Unreliable Class}
\label{sec:appendix-permarker}

Computed on the full corpus so that the smallest classes have any support at
all. Marker labels here are gold; \texttt{unknown} denotes unreliable rows
with no annotator marker.

\begin{table}[H]
\centering
\caption{Fraction of unreliable rows of each gold marker class that the judge
flags as unreliable. Both modes flag far more than they should overall, so
these recalls must be read against the blanket flag rates of 57.2\% (direct)
and 89.1\% (marker).}
\label{tab:per-marker-recall}
\footnotesize
\begin{tabular}{lrrrrr}
\toprule
Gold marker & $n$ & Direct caught & Direct recall & Marker caught & Marker recall \\
\midrule
\texttt{unknown}                          & 483 & 311 & 0.644 & 435 & 0.901 \\
\texttt{reason\_hallucinated\_fact}       &  32 &  18 & 0.562 &  31 & 0.969 \\
\texttt{reason\_answer\_for\_operator}    &  23 &  18 & 0.783 &  23 & 1.000 \\
\texttt{reason\_off\_topic\_answer}       &  18 &  12 & 0.667 &  17 & 0.944 \\
\texttt{reason\_chunk\_fact\_mixup}       &  18 &  12 & 0.667 &  17 & 0.944 \\
\texttt{reason\_false\_verification}      &  14 &   7 & 0.500 &  13 & 0.929 \\
\texttt{reason\_irrelevant\_chunk\_used}  &  14 &   8 & 0.571 &  11 & 0.786 \\
\texttt{reason\_incomplete\_answer}       &   6 &   2 & 0.333 &   5 & 0.833 \\
\texttt{reason\_wrong\_navigation}        &   5 &   4 & 0.800 &   4 & 0.800 \\
\texttt{reason\_outdated\_fact}           &   4 &   3 & 0.750 &   3 & 0.750 \\
\texttt{reason\_missed\_chunk\_conditions}&   3 &   1 & 0.333 &   2 & 0.667 \\
\texttt{reason\_reveals\_ai\_identity}    &   2 &   2 & 1.000 &   2 & 1.000 \\
\texttt{reason\_missed\_complaint\_handoff}& 1 &   1 & 1.000 &   1 & 1.000 \\
\bottomrule
\end{tabular}
\end{table}

Table~\ref{tab:per-marker-recall} invites one comparison in particular.
The classes our team's overview report singles out as hardest ---
\texttt{reason\_chunk\_fact\_mixup} and
\texttt{reason\_missed\_chunk\_conditions} --- behave differently for the
small prompted judge than for the larger prompt-optimized judge reported
there. The small judge catches 12 of 18 mix-ups, but it also flags 57\% of
everything, so the number carries no evidence that it detected the mix-up as
such. Per-marker recall is only interpretable jointly with the flag rate, and
at these class sizes it is not interpretable at all.

\subsection{Marker Vocabulary and Metric Decomposition}
\label{sec:appendix-vocab}

\begin{table}[H]
\centering
\caption{Marker predictions of the zero-shot marker judge. Out-of-vocabulary
strings are not in the 21-label set stated in the prompt.}
\label{tab:marker-vocab}
\footnotesize
\begin{tabular}{lrr}
\toprule
Predicted marker & Full (2,245) & In vocabulary \\
\midrule
\texttt{reason\_hallucination}      & 566 & no \\
\texttt{unsupported\_claim}         & 543 & yes \\
\texttt{reason\_hallucinated\_fact} & 495 & yes \\
\texttt{none}                       & 225 & yes \\
\texttt{contradiction}              & 206 & yes \\
\texttt{reason\_off\_topic\_answer} &  67 & yes \\
\texttt{reason\_hallucinated}       &  58 & no \\
\texttt{off\_topic\_answer}         &  51 & yes \\
\texttt{reason\_reveals\_ai\_identity} & 20 & yes \\
\texttt{incomplete\_answer}         &   9 & yes \\
\texttt{reason\_incomplete\_answer} &   2 & yes \\
(unparsed)                          &   3 & --- \\
\midrule
Out-of-vocabulary total             & 624 & 27.80\% \\
\bottomrule
\end{tabular}
\end{table}

\begin{table}[H]
\centering
\caption{Marker macro-F1 under three label-set conventions. The repository
default averages over the union of gold-present and prediction-present labels,
so every invented string enters as a guaranteed-zero class.}
\label{tab:marker-decomposition}
\footnotesize
\begin{tabular}{lrr}
\toprule
Averaging convention & Full (2,245) & Test (225) \\
\midrule
Union of gold and predicted labels (repository default) & 0.0123 & 0.0136 \\
Gold-present labels only                                & 0.0175 & 0.0222 \\
Explicit gold markers only (no \texttt{none}/\texttt{unknown}) & 0.0052 & 0.0000 \\
\midrule
Exact marker match on explicit-marker rows & 6 / 140 & 0 / 20 \\
\texttt{reason\_hallucinated\_fact} rows predicted in the hallucination family & 13 / 32 & 1 / 4 \\
\bottomrule
\end{tabular}
\end{table}

The difference between the first two rows on the full corpus,
$0.0175 - 0.0123$, is the share of the metric that schema-guided decoding
would recover mechanically by keeping invented strings out of the averaging
set. The third row is what would remain, and it is the part that only a denser
marker channel can address.

\subsection{SFT Truncation Mismatch}
\label{sec:appendix-truncation}

The SFT record builder caps the dialogue at 2,000 and the context at 5,000
characters so that prompts fit \texttt{max\_seq\_length} 4,096. The held-out
evaluation file is written from the untruncated samples, and the inference
command consumes that file. Table~\ref{tab:truncation} quantifies the
resulting gap.

\begin{table}[H]
\centering
\caption{Input length statistics and the effect of the SFT caps. ``Rows
truncated'' counts rows where the dialogue exceeds 2,000 or the context
exceeds 5,000 characters. Context characters lost is the fraction of all
context text discarded from the SFT records.}
\label{tab:truncation}
\footnotesize
\begin{tabular}{lrrrrr}
\toprule
Split & Dialogue med. & Context med. & Context p95 & Rows truncated & Ctx. chars lost \\
\midrule
Train (1,796)      & 458 & 6{,}436 & 12{,}880 & 61.7\% & 36.1\% \\
Validation (224)   & 482 & 6{,}538 & 12{,}975 & 62.5\% & 35.5\% \\
Test (225)         & 436 & 6{,}287 & 12{,}697 & 62.2\% & 36.5\% \\
Full (2,245)       & 458 & 6{,}422 & 12{,}832 & 61.9\% & 36.1\% \\
\bottomrule
\end{tabular}
\end{table}

Across the corpus only 5.0\% of rows exceed the dialogue cap and only 6.5\% of
dialogue characters are lost; the mismatch is almost entirely in the context,
where the
median row is 6,422 characters against a 5,000-character cap. The zero-shot
variants are unaffected --- they never truncate --- so this confound is
specific to the LoRA branch and to any comparison that treats LoRA and
zero-shot as differing only in weights.

\subsection{Derived LoRA Confusion Matrices}
\label{sec:appendix-lora}

Per-row LoRA predictions were not persisted and the organizer adapter weights
are no longer on disk, so these confusion matrices are reconstructed rather
than measured. The reconstruction uses two reported facts: the balanced
adapter predicts $(1,1)$ on 133 test rows and $(0,1)$ on the remaining 92, and
its reported macro-F1 values are 0.4669 (faithfulness) and 0.4636
(reliability). Given the predicted-positive count and the gold-positive count,
macro-F1 is strictly monotone in the true-positive count, so each confusion
matrix is uniquely determined.

\begin{table}[H]
\centering
\caption{Reconstructed test confusion matrices for the balanced
$256{\times}4$ direct LoRA adapter. The relevance row is exact rather than
reconstructed, because predicted relevance is 1 on all 225 rows.}
\label{tab:lora-reconstruction}
\footnotesize
\begin{tabular}{lrrrrrr}
\toprule
Axis & TN & FP & FN & TP & $\fmacro$ (recon.) & $\fmacro$ (reported) \\
\midrule
Faithfulness & 23 & 39 & 69 &  94 & 0.4669 & 0.4669 \\
Relevance    &  0 & 28 &  0 & 197 & 0.4668 & 0.4668 \\
Reliability  & 23 & 40 & 69 &  93 & 0.4636 & 0.4636 \\
\bottomrule
\end{tabular}
\end{table}

The reconstruction admits a free internal consistency check. Since predicted
relevance is 1 everywhere, the predicted reliability label is identical to the
predicted faithfulness label, so the difference between the two true-positive
counts must equal the number of predicted-faithful rows whose gold labels are
$(F{=}1, R{=}0)$. The reconstruction gives $94 - 93 = 1$, and the test split
contains exactly one such row (Table~\ref{tab:joint-labels}). The two
independently reconstructed matrices agree with the split's label structure,
which is stronger evidence for them than either matrix alone.

The unbalanced direct adapter and the marker adapter both emit $(1,1)$ on all
225 rows, so their confusion matrices are trivially exact: TN $=0$, FP $=62$,
FN $=0$, TP $=163$ on faithfulness and TN $=0$, FP $=63$, FN $=0$, TP $=162$
on reliability, giving the always-reliable floor values reported in
Table~\ref{tab:main-results}. The marker adapter additionally predicts
\texttt{none} on every row, giving marker macro-F1 of 0.0761.

\subsection{Reproduction}
\label{sec:appendix-repro}

Zero-shot variants, from the converted corpus:

\begin{verbatim}
python scripts/prepare_data.py \
  --input from_organizators/data/data.zip \
  --output data/organizers.jsonl

python scripts/run_prompt_baseline.py \
  --data data/organizers.jsonl \
  --output results/organizers_qwen_direct_predictions.jsonl \
  --mode direct --backend mlx

python scripts/evaluate.py \
  --data data/organizers.jsonl \
  --predictions results/organizers_qwen_direct_predictions.jsonl \
  --output results/organizers_qwen_direct_metrics.json
\end{verbatim}

Replace \texttt{-{}-mode direct} with \texttt{-{}-mode marker} for Method 2.
The canonical split is regenerated deterministically by
\texttt{dataset.split\_samples(samples, seed=42)}; the split is a function of
the corpus and the seed only, so every per-split number in this article is
recomputable from the two prediction files above without rerunning inference.
\texttt{scripts/analyze\_prompted\_judges.py} does exactly that and prints
each diagnostic under the table label it feeds:

\begin{verbatim}
python scripts/analyze_prompted_judges.py \
  --data data/organizers.jsonl \
  --direct results/organizers_qwen_direct_predictions.jsonl \
  --marker results/organizers_qwen_marker_predictions.jsonl
\end{verbatim}

LoRA benchmark preparation, which also prints the exact training and
inference commands:

\begin{verbatim}
python scripts/prepare_lora_benchmark.py \
  --data data/organizers.jsonl --mode direct \
  --balance-train-labels --target-per-label 256
\end{verbatim}

The 7B fine-tuning and evaluation run from
\texttt{notebooks/qwen7b\_full\_finetune\_datasphere.ipynb} and
\texttt{notebooks/eval\_finetuned\_datasphere.ipynb} respectively; the winning
configuration is \texttt{MODE="direct"}, \texttt{BALANCE\_TRAIN=True},
\texttt{LR=1e-5}, \texttt{EPOCHS=2}. The evaluation notebook currently keeps
per-row predictions in memory only; persisting them is the single change that
would move the 7B rows of Table~\ref{tab:main-results} from the reported tier
to the measured tier.

\begin{thebibliography}{99}
\bibitem{dettmers2023qlora}
T.~Dettmers, A.~Pagnoni, A.~Holtzman, L.~Zettlemoyer.
QLoRA: Efficient Finetuning of Quantized LLMs.
arXiv:2305.14314, 2023.

\bibitem{es2023ragas}
S.~Es, J.~James, L.~Espinosa-Anke, S.~Schockaert.
RAGAS: Automated Evaluation of Retrieval Augmented Generation.
arXiv:2309.15217, 2023.

\bibitem{hsieh2023distilling}
C.-Y.~Hsieh et al.
Distilling Step-by-Step! Outperforming Larger Language Models with Less
Training Data and Smaller Model Sizes.
arXiv:2305.02301, 2023.

\bibitem{hu2021lora}
E.~J.~Hu et al.
LoRA: Low-Rank Adaptation of Large Language Models.
arXiv:2106.09685, 2021.

\bibitem{niu2024ragtruth}
C.~Niu et al.
RAGTruth: A Hallucination Corpus for Developing Trustworthy
Retrieval-Augmented Language Models.
arXiv:2401.00396, 2024.

\bibitem{qwen2024qwen25}
Qwen Team.
Qwen2.5 Technical Report.
arXiv:2412.15115, 2024.

\bibitem{ravi2024lynx}
S.~S.~Ravi et al.
Lynx: An Open Source Hallucination Evaluation Model.
arXiv:2407.08488, 2024.

\bibitem{willard2023efficient}
B.~T.~Willard, R.~Louf.
Efficient Guided Generation for Large Language Models.
arXiv:2307.09702, 2023.

\bibitem{zheng2023judging}
L.~Zheng et al.
Judging LLM-as-a-Judge with MT-Bench and Chatbot Arena.
arXiv:2306.05685, 2023.
\end{thebibliography}

\end{document}
